\chapter{Modern Portfolio Theory}
% See https://en.wikipedia.org/wiki/Outline_of_finance#Portfolio_theory
% Also Bessel's correction for sample means and deviations.
% https://en.wikipedia.org/wiki/Bessel%27s_correction

\section{Cross-Correlation}
In signal processing, {\bf cross-correlation}\footnote{Cross-correlation is
also called the {\bf sliding dot product} or {\bf sliding inner product}}
is a measure of similarity of two signals.  

In probability theory, cross-correlation refers to the correlations between
two random variables $X$ and $Y$.


\section{Portfolio and Assets}

Let a portfolio $P$ be a collection of $n$ individual assets 
with present values $a_1, a_2, \ldots a_n$.
Assets have units of 
U.S.~dollars (USD).\footnote{When the asset is stock, bond, or 
other instrument that consists of 
fractional ownership (shares) denominated by price per share, 
then the present value of the asset is simply the number of shares
held (holdings) in the portfolio multiplied by the price per share of that
holding.}

The present valuation of the portfolio, $PV(P)$, is given as
\be
 PV(P) = \sum_{i=1}^n a_i.
\ee

\section{Portfolio Weighting}

The relative weight $w_i$ of each asset $a_i$ is simply the fraction 
of the asset's present value to the portfolio's present value, thus
\be
 w_i = \frac{a_i}{PV} = \frac{a_i}{\sum_{i=1}^n a_i}.
\ee
The sum of all weights, must equal one (1).
\be
 \sum_{i=1}^n w_i = \sum_{i=1}^n \frac{a_i}{PV} = \frac{1}{PV} \sum_{i=1}^n a_i 
 = \frac{PV}{PV} = 1.
\ee

\section{Portfolio Return}

Let $\E(R_P)$ and $\E(R_i)$ the be expected periodic return $R_P$ 
of the portfolio $P$ and 
expected periodic return $R_i$ on an asset $a_i$, 
respectively.\footnote{The period
is typically stated annually, quarterly, or yearly.}
The expected portfolio return is then simply the weighted sum of each 
asset's return, 
\be
\E(R_P) = \sum_{i=1}^n w_i \; \E(R_i).
\ee

Let the mean periodic return $\bar{R}(a_i)$ of asset $a_i$ is
given by observing the return at $K$ times during the periodic return 
interval, 
\be
\bar{R}(a_i) = \frac{1}{K} \sum_{k=1}^K R_k(a_i).
\ee
For clarity, the notation $f(a_i)$ to denote {\em on a per-asset basis} 
is assumed, as written hereinafter simply as $f$.
Then the mean periodic return is given as
\be
 \bar{R} = \frac{1}{K} \sum_{k=1}^K R_k.
\ee
Let $\var(R) = s^2$ be the (sample) variance of periodic returns on an 
asset $a_i$, given by
\be
 \var(R) = s^2 = \frac{1}{K-1} \sum_{k=1}^K \left(R_k - \bar{R}\right)^2.
\ee
The (sample) standard deviation $s$ is the positive square root of the 
(sample) variance $s^2$,
\be
 s = \sqrt{\frac{1}{K-1} \sum_{k=1}^K \left(R_k - \bar{R}\right)^2}.
\ee

